% ======================================================================
% SECTION 3: HR 9503 - Cancer Patient Procedural Modification
% Authority Act of 2026
% Adaption of OHRP Broad Consent 2017 and 21st Century Cures Act 2016.
%
% This file consolidates instructions previously housed in the prior
% template's Patient Priority Framework dimension 2.2 (enrollment
% comprehension), Implementation/Metrics (Domain 2 decision quality),
% Discussion (11.1.a Cures Act extension and 11.3.d OHRP Broad
% Consent opening), Limitations (12.1.c HR 4503 limitation), and
% Conclusions (13.1 third list item) into the body of HR 9503.
% ======================================================================

[CONSOLIDATED PATIENT PRIORITY FRAMEWORK DIMENSION 2.2 (formerly
Patient Priority Framework subsection 2.2). Enrollment comprehension:
key information rules, broad consent, patient-initiated translation
of plain-language summaries, real-time consent updates. Cite
\cite{fda_informed_consent_guidance_2023},
\cite{fda_ohrp_key_information_draft_2024}, and
\cite{ohrp_broad_consent_attachment_c_2017}.]

[CONSOLIDATED INTRODUCTION TIE-IN (formerly Introduction block 2 item
2.g). FDA Informed Consent Guidance 2023, FDA/OHRP Key Information
Draft 2024, and OHRP Broad Consent Recommendations 2017 establish the
prior consent and broad-consent frameworks. HR 9503 closes the
patient-control gap by adding a positive procedural-modification
right where the patient may modify their own treatment under
physical-AI execution. Cite \cite{fda_informed_consent_guidance_2023},
\cite{fda_ohrp_key_information_draft_2024}, and
\cite{ohrp_broad_consent_attachment_c_2017}.]

[PROCESSING INSTRUCTION (HR 9503 - Cancer Patient Procedural
Modification Authority Act of 2026, populate to 9 to 12 paragraphs
structured into the five legislative-act subsections):

The bill must explicitly state in its first paragraph that it is an
"Adaption of OHRP Broad Consent Recommendations 2017 and 21st Century
Cures Act of 2016 (PL 114-255)" so that readers can trace the chain
of prior authority. Cite \cite{ohrp_broad_consent_attachment_c_2017},
\cite{rank1_cures_act_2016}, and
\cite{hhs_45cfr46_protection_human_subjects} for 45 CFR 46 (Common
Rule) baseline.

Subsection 3.1 - Findings and Declarations: Read
patients/paper/Deep-Research-3/part1_legal_baseline.md (the
broad-consent permissions for secondary biospecimen use) and
patient-journey/stage_07_immunotherapy.py (the 35 pembrolizumab
cycles managed by Franka Emika at USL 88.75) to establish that
current procedure consent is binary (accept entire protocol or
refuse) and that physical-AI execution now permits per-cycle and
per-step patient modification with no loss of safety. Compose 5 to 7
enumerated findings citing \cite{kawchak_2026_19119939},
\cite{kawchak_2026_18973368} (the ICH E6(R3) adaption that already
permits such modifications under "real-time consent updates"), and
\cite{kawchak_2026_19040707} (the 21 CFR Part 50 adaption).

Subsection 3.2 - Definitions:
- "Procedural modification" means a patient-initiated change to a
  pre-approved protocol step, executed by the qualified Physical AI
  system, that does NOT alter the primary endpoint or the safety
  profile beyond pre-defined tolerances. Reference the
  ICH E6(R3) adaption at
  national-platform/ich_e6r3_adapt/01_preamble_principles_investigator.tex
  and 02_sponsor_responsibilities.tex (read both files).
- "Real-time consent update" means a HIPAA-compliant, FHIR-based
  consent revision logged in the patient's EHR within 5 minutes of
  the modification. Reference the FHIR infrastructure at
  national-platform/national_mcp/chunk2_methods_and_results_part1.tex
  and chunk3_results_part2.tex (the MCP-PAI five-server topology with
  hash-chained audit trails).
- "Qualified Physical AI executor" means a robot or robotic system
  with a USL score above the modification-tier threshold (e.g.,
  USL 8.0 or above for surgical modifications, 6.0 or above for
  therapeutic positioning modifications, 4.0 or above for
  rehabilitation modifications). Cite \cite{kawchak_2026_18778219}.

Subsection 3.3 - Patient Procedural Modification Rights (operative
core):
- Right to request, at any point during the trial, a modification to
  a pre-approved procedural step, with the qualified Physical AI
  executor evaluating safety in real time and either executing or
  declining the modification with cited reason.
- Right to receive, in plain language, the safety evaluation result
  for any requested modification within 5 minutes; if declined, the
  patient may appeal to the IRB on call (the
  21cfr50_adapt/02_irb_review_pediatric.tex IRB framework applies).
- Right to revise broad consent at any point with 24-hour notice,
  drawing on patients/paper/Deep-Research-3/part1_legal_baseline.md
  broad-consent baseline and \cite{ohrp_broad_consent_attachment_c_2017}.
- Right to specify, before each procedure, which steps may be
  modified by the patient and which require investigator review.
  Reference the 10-stage patient journey in patient-journey/ and the
  per-stage modification points enumerated at
  patient-journey/stage_05_surgery.py through
  stage_10_closeout.py.
- Right to request that a procedural step be performed by a
  Physical AI system rather than by a human doctor or nurse where
  doing so reduces the documented error rate (this dovetails with
  HR 9504 in the next section).
- Right to a per-cycle "modifiable steps" matrix delivered as part of
  the trial enrollment package, listing every step that the patient
  may modify under broad consent without requiring fresh IRB review.

Subsection 3.4 - Implementation: 18-month deadline for federal IRB
networks to publish a per-protocol "modifiable steps" matrix; FDA,
OHRP, and HHS roles; reference the prior site documentation at
new-trial/site/02-legislation-patient-rights/physical_ai_patient_rights.tex
and at
national-platform/new_trial_psl/02_ab2847_patient_rights.tex for the
California baseline.

Subsection 3.5 - Reporting and Enforcement: Annual federal report on
modification rate by stage, by cancer, by robot category; the audit
trail manager at sponsor/final_paper/scripts/core_agents/
audit_trail_manager.py provides the technical baseline; civil
penalties harmonized with 45 CFR 46 enforcement and the Cures Act
information-blocking framework. Cite
\cite{federal_register_info_blocking_disincentives_2024}.]

[CONSOLIDATED IMPLEMENTATION-METRICS DOMAIN 2 - DECISION QUALITY
(formerly Implementation Metrics subsection 10.2 Domain 2). The
federal Patient Control Dashboard must publish, for HR 9503, the
following Domain 2 metrics: consent comprehension, decisional-
conflict scores, goal-trial concordance, AI-exclusion appeal success
rate, explanation and provenance coverage. Tied to the HTI-1 DSI
transparency requirements at \cite{hti1_dsi_fact_sheet_2023}.]

[CONSOLIDATED DISCUSSION COMPARISON (formerly Discussion items
11.1.a Cures Act extension and 11.3.d OHRP Broad Consent opening).
HR 9503 extends the prior 21st Century Cures Act broad-consent
permissions into a positive procedural modification authority and
extends the OHRP Broad Consent Recommendations 2017 into a
real-time consent update framework. Cite
\cite{rank1_cures_act_2016},
\cite{onc_cures_act_final_rule_page}, and
\cite{ohrp_broad_consent_attachment_c_2017}.]

[CONSOLIDATED LIMITATIONS PARAGRAPH (formerly Limitations item
12.1.c). HR 9503 assumes that real-time consent updates can be
logged within 5 minutes by every qualified Physical AI executor.
The MCP-PAI infrastructure is documented; field deployment is not.
The bill's "modifiable steps" matrix presumes IRB review at scale
that has not been tested. Cite \cite{kawchak_2026_18973368}.]

[CONSOLIDATED CONCLUSIONS RESTATEMENT (formerly Conclusions item
13.1 third list item). HR 9503 replaces the prior all-or-nothing
consent pattern under traditional IRB review with a granular
per-step national standard for patient-led procedural modification.
Cite \cite{kawchak_2026_18973368} and \cite{kawchak_2026_19119939}.]

Close the bill (1 paragraph after Subsection 3.5) with two patient-
control framings:

(i) FOR INDIVIDUAL PATIENTS: A 58-year-old female NSCLC patient
    matching PAT-2026-0042 (patient-journey/stage_07_immunotherapy.py)
    requests at cycle 22 of 35 pembrolizumab that the Franka Emika
    cobot deliver the infusion 30 minutes earlier than the protocol
    default to accommodate work, with the qualified Physical AI
    executor evaluating safety in 4 minutes and approving the
    modification under the patient's broad-consent revision.

(ii) FOR BROAD ADOPTION: All Physical AI oncology trial sites must
     publish a "modifiable steps" matrix per protocol within 18
     months. The bill positions the United States as Number 1 in
     the world for patient-led procedural modification under
     physical-AI execution, replacing the prior all-or-nothing
     consent pattern under traditional IRB review with a granular
     per-step national standard. Cite \cite{kawchak_2026_18973368}
     and \cite{kawchak_2026_19119939}.
